\documentclass[UTF8,fontset=none,aspectratio=169,10pt]{beamer}
\usepackage{ads-beamer-theme}
\title{高级分布式系统}
\subtitle{第十四讲：Consensus、Raft 与一致性日志}
\author{陈鹏飞\quad 计算机学院}
\institute{中山大学}
\date{2025--2026 学年}

\begin{document}
\begin{frame}[plain]
\begin{tikzpicture}[remember picture,overlay]
  \fill[ADSBlue] (current page.north west) rectangle ([xshift=2.4cm]current page.south west);
  \fill[ADSOrange] ([xshift=2.4cm]current page.north west) rectangle ([xshift=2.55cm]current page.south west);
  \node[anchor=west,align=left] at ([xshift=3.3cm,yshift=1.0cm]current page.west) {\fontsize{32}{38}\selectfont\color{ADSBlue}\bfseries 高级分布式系统};
  \node[anchor=west,align=left] at ([xshift=3.3cm,yshift=-0.4cm]current page.west) {\fontsize{21}{27}\selectfont\color{ADSDarkBlue}第十四讲\quad Consensus 与一致性日志};
  \node[anchor=west,align=left] at ([xshift=3.3cm,yshift=-2.2cm]current page.west) {\large\color{ADSGray}让所有正确副本对同一个历史达成一致};
  \node[anchor=east] at ([xshift=-0.8cm,yshift=-3.2cm]current page.east) {\small\color{ADSGray}陈鹏飞\quad 中山大学\quad 2025--2026};
  \draw[ADSBlue,line width=0.8pt] ([xshift=3.3cm,yshift=-1.35cm]current page.west) -- ([xshift=12.6cm,yshift=-1.35cm]current page.west);
\end{tikzpicture}
\end{frame}

\sectionframe{一、共识到底要解决什么？}{不是“大家都收到消息”，而是所有正确进程最终选择同一个、可接受的值或日志前缀。}

\begin{frame}{本讲学习目标}
\begin{enumerate}
  \item 区分 Byzantine Agreement、Consensus 和 Interactive Consistency；
  \item 掌握 Termination、Agreement、Integrity、Validity；
  \item 解释同步/异步、Crash/Byzantine 对可解性的影响；
  \item 理解多数派、任期、日志匹配和提交索引；
  \item 把共识映射到数据库、Kubernetes 控制面和 AI 模型注册。
\end{enumerate}
\vfill
\callout{主线}{共识的对象可以是一个值，也可以是一串经过排序、复制和提交的日志命令。}
\end{frame}

\begin{frame}{三种相关问题}
\begin{tabularx}{\linewidth}{>{\bfseries}p{3.2cm}XX}
\toprule
Byzantine Agreement & 一个 source 提出初值；正确进程一致决定，source 正确时保持其值。 \\
Consensus & 所有进程都有初值；若所有正确进程初值相同，决定必须保持该值。 \\
Interactive Consistency & 所有正确进程对一整个值数组达成一致，故障进程位置可任意填充。 \\
\bottomrule
\end{tabularx}
\vfill
\callout{共同结构}{它们可以通过规约互相构造，核心都是“同一决定 + 可接受的有效性”。}
\end{frame}

\begin{frame}{共识的四个性质}
\begin{enumerate}
  \item Termination：每个正确进程最终决定；
  \item Agreement：任意两个正确进程不能决定不同值；
  \item Integrity：每个进程最多决定一次；
  \item Validity：决定值必须来自某个提议，不能凭空创造。
\end{enumerate}
\vfill
\callout{为什么要分开}{Agreement 是安全性，Termination 是活性；Failure Detector 主要帮助解决活性，但不能牺牲安全性。}
\end{frame}

\sectionframe{二、模型决定上限}{先写清时钟、网络和故障模型，再讨论算法是否可行。}

\begin{frame}{故障模型与可解性}
\begin{tabularx}{\linewidth}{>{\bfseries}p{2.6cm}XX}
\toprule
无故障 & 同步/异步都可达成一致；直接交换值即可。 \\
Crash & 同步消息传递可容忍 $f<n$；异步系统需额外 Failure Detector。 \\
Byzantine & 同步系统在 $f\le\lfloor(n-1)/3\rfloor$ 时可达成；异步 Byzantine 更困难。 \\
\bottomrule
\end{tabularx}
\vfill
\callout{直觉}{故障越强、时间信息越弱，系统越需要冗余和更强的交互轮次。}
\end{frame}

\begin{frame}{同步 Crash 共识：f+1 轮为何足够？}
\begin{enumerate}
  \item 最多允许 $f$ 个进程崩溃；
  \item 执行 $f+1$ 轮广播，每轮传播当前最小/最高优先级估计；
  \item 至少有一轮没有新的崩溃；
  \item 该轮之后所有正确进程能看到相同的候选值；
  \item 输出相同值满足 Agreement，值来自提议满足 Validity。
\end{enumerate}
\end{frame}

\begin{frame}{异步系统中的 FLP 直觉}
\begin{itemize}
  \item 如果一个进程迟迟没有消息，其他进程无法知道它是慢还是崩溃；
  \item 立即决定可能在迟到消息到达后产生分歧；
  \item 无限等待则无法满足 Termination；
  \item 因此纯异步、允许一个 crash failure 时，不存在同时满足所有性质的确定性共识算法。
\end{itemize}
\vfill
\callout{工程策略}{引入最终准确 Failure Detector、随机化、同步窗口或外部共识服务。}
\end{frame}

\sectionframe{三、从多数派到日志复制}{Raft/Paxos 的共同骨架：任期、leader、多数派和不可回退的日志前缀。}

\begin{frame}{多数派交集为什么重要？}
\begin{itemize}
  \item 集群有 $N$ 个副本，提交通常需要多数派 $\lfloor N/2\rfloor+1$；
  \item 任意两个多数派至少有一个交集副本；
  \item 已提交的命令不能在另一条多数派路径上被替换；
  \item 多数派同时提供容错：最多容忍少于半数副本失效。
\end{itemize}
\end{frame}

\begin{frame}{Raft 的状态机}
\begin{center}
\begin{tikzpicture}[node distance=9mm and 13mm,>=Latex]
  \nodebox[minimum width=2.4cm]{ADSBlue}{follower}{Follower};
  \nodebox[right=of follower,minimum width=2.4cm]{ADSOrange}{candidate}{Candidate};
  \nodebox[right=of candidate,minimum width=2.4cm]{ADSGreen}{leader}{Leader};
  \draw[\arrow,ADSGray] (follower)--node[above,font=\scriptsize]{选举超时}(candidate);
  \draw[\arrow,ADSGray] (candidate)--node[above,font=\scriptsize]{多数派票}(leader);
  \draw[->,ADSGray] (leader.south) .. controls +(0,-1) and +(0,-1) .. node[below,font=\scriptsize]{更高任期/失去心跳}(follower.south);
\end{tikzpicture}
\end{center}
\begin{itemize}
  \item 任期 term 让旧 leader 的消息可被识别为过期；
  \item 同一任期最多一个合法 leader，是安全性的关键。
\end{itemize}
\end{frame}

\begin{frame}{Leader Election 的因果链}
\begin{enumerate}
  \item Follower 在 election timeout 内没有收到心跳；
  \item 增加 term，转换为 Candidate，投自己一票；
  \item 向其他副本发送 RequestVote；
  \item 获得多数派后成为 Leader，立即发送心跳；
  \item 若看到更高 term，立刻退回 Follower。
\end{enumerate}
\end{frame}

\begin{frame}{日志复制：AppendEntries}
\begin{itemize}
  \item Leader 把客户端命令追加到本地日志；
  \item 向 Followers 发送 prevLogIndex、prevLogTerm 和新 entries；
  \item Follower 只有在前缀匹配时接受新 entries；
  \item Leader 收到多数派确认后推进 commitIndex；
  \item 各副本按 commitIndex 将日志应用到状态机。
\end{itemize}
\end{frame}

\begin{frame}{日志匹配性质}
\begin{itemize}
  \item 如果两个日志在相同 index、term 相同，则该 index 之前的日志也相同；
  \item Leader 只覆盖冲突的未提交后缀，不覆盖已提交前缀；
  \item 旧 Leader 恢复后通过 term 比较和日志回退变成 Follower；
  \item 状态机只执行 committed entries，因此不会执行分叉历史。
\end{itemize}
\end{frame}

\begin{frame}{提交与线性化}
\begin{itemize}
  \item 多数派确认意味着命令即使丢失少数副本也能恢复；
  \item 客户端收到成功响应通常要等到 entry committed；
  \item 线性化还需要处理客户端重试、重复请求和 leader 变化；
  \item 常见做法是 client request ID + 去重表 + read barrier。
\end{itemize}
\end{frame}

\sectionframe{四、Paxos/共识的证明视角}{不同算法名字背后，都是“不会出现两个互相冲突的已选值”。}

\begin{frame}{Paxos 的角色抽象}
\begin{itemize}
  \item Proposer：提出候选值；
  \item Acceptor：承诺、接受提案；
  \item Learner：学习已经选定的值；
  \item 提案编号全序化，后续提案必须尊重已被多数派接受的值；
  \item 两阶段 prepare/accept 让新 proposer 了解旧历史。
\end{itemize}
\end{frame}

\begin{frame}{Paxos 的不变量}
\begin{enumerate}
  \item 一个提案编号最多对应一个值；
  \item 若值 v 已被某多数派接受，之后更高编号的多数派交集必能看到 v；
  \item 新 proposer 不能随意选择与已选值冲突的值；
  \item 因而不可能出现两个不同的已选值。
\end{enumerate}
\vfill
\callout{理解方法}{先画两个多数派的交集，再问交集里的 acceptor 带来了什么历史信息。}
\end{frame}

\begin{frame}{Paxos 完整例子：后来者为什么不能改值}
\small
三个 acceptor 为 $A,B,C$，任意两个构成多数派。
\begin{center}
\begin{tabularx}{.96\linewidth}{>{\bfseries}p{1.2cm}p{4.2cm}X}
\toprule
阶段 & 消息 & 状态变化 \\
\midrule
1 & P1: PREPARE(5) $\rightarrow$ A,B & A、B promise 不再接受编号 $<5$ 的提案。 \\
2 & P1: ACCEPT(5,X) $\rightarrow$ A,B & A、B 接受 $(5,X)$；$X$ 已被多数派选定。 \\
3 & P2: PREPARE(8) $\rightarrow$ B,C & B 回复已接受 $(5,X)$；C 回复“无已接受值”。 \\
4 & P2 选择值 & P2 必须采用回复中最高编号的已接受值 $X$，不能改提 $Y$。 \\
5 & P2: ACCEPT(8,X) $\rightarrow$ B,C & 新多数派接受的仍是 $X$。 \\
\bottomrule
\end{tabularx}
\end{center}
\callout{安全性来源}{两个多数派 $\{A,B\}$ 与 $\{B,C\}$ 在 B 相交；B 把已选历史带给更高编号提案。}
\end{frame}

\begin{frame}{Raft 与 Paxos 的教学对比}
\begin{tabularx}{\linewidth}{>{\bfseries}p{2.7cm}XX}
\toprule
领导者 & Raft 显式维护一个 leader；Paxos 可有多个 proposer，但通常工程实现也选稳定 leader。 \\
日志 & Raft 把日志复制写成完整协议；Paxos 需要 Multi-Paxos 组织连续命令。 \\
可理解性 & Raft 的角色、任期和状态机更易讲授；Paxos 的不变量更适合证明。 \\
\bottomrule
\end{tabularx}
\end{frame}

\sectionframe{五、共识在工程系统中的边界}{共识保证副本对同一顺序达成一致，但不自动解决业务幂等、外部副作用和性能。}

\begin{frame}{Kubernetes 控制面的共识需求}
\begin{itemize}
  \item etcd 用 Raft 一致性日志保存对象、租约和 watch 所需状态；具体读路径仍需区分线性化读与可能的串行化/缓存读；
  \item API Server 可以水平扩展，但写入顺序由共识存储维护；
  \item 控制器依赖线性化读取或版本检查，避免基于旧对象重复操作；
  \item etcd 不可用时，旧状态可能仍被缓存，但新的期望状态不能安全提交。
\end{itemize}
\end{frame}

\begin{frame}{近期案例：etcd/Raft 如何支撑控制面}
\begin{enumerate}
  \item API Server 把对象写请求提交给 etcd；Raft 先让多数派副本对日志顺序达成一致；
  \item 控制器通过 watch 观察版本推进，再把实际状态写回新的对象版本；
  \item 网络分区时，少数派可以保留旧缓存，但不能安全地产生新的线性化写入；
  \item 因此“控制器最终收敛”依赖两层机制：共识日志提供稳定事实，reconcile loop 处理事实与现实之间的差距。
\end{enumerate}
\vfill
\callout{分析题}{如果 watch 连接断开后重新建立，为什么必须用 resource version 做重列举（relist），而不能只相信断线前收到的最后一条事件？}
\end{frame}

\begin{frame}{AI 模型注册与发布}
\begin{enumerate}
  \item 多个训练/评测流程产生候选模型；
  \item 注册中心用共识日志记录 model version、指标和审批状态；
  \item 只有一个版本能成为 production alias；
  \item 推理控制器 watch 到提交事件后滚动升级；
  \item 旧版本、灰度比例和回滚命令都要可审计、可重放。
\end{enumerate}
\end{frame}

\begin{frame}{共识不能替代的东西}
\begin{itemize}
  \item 共识不能让外部支付、邮件或 GPU 实际操作自动回滚；
  \item 外部副作用需要幂等键、事务 outbox、补偿或 fencing；
  \item 共识日志只保证命令顺序，不保证命令执行时间短；
  \item 大对象（模型权重、数据集）不应直接写入共识日志，应只记录不可变引用。
\end{itemize}
\end{frame}

\sectionframe{六、课堂推演与 AI 练习}{用任期、日志和多数派把“谁能决定”变成可执行的状态机。}

\begin{frame}{课堂推演：旧 Leader 恢复}
\begin{enumerate}
  \item term=5 的 Leader A 网络分区；
  \item B、C、D 选出 term=6 的新 Leader B 并提交命令 x；
  \item A 恢复，仍向客户端返回 term=5 的旧日志；
  \item 问题：A 如何发现自己过期？旧日志怎样被截断？客户端怎样避免重复执行？
\end{enumerate}
\end{frame}

\begin{frame}{课堂推演：多数派与故障}
\begin{tabularx}{\linewidth}{>{\bfseries}p{2.8cm}XX}
\toprule
副本数 & 可容忍故障 & 提交条件 \\
3 & 1 & 2 个副本确认 \\
5 & 2 & 3 个副本确认 \\
7 & 3 & 4 个副本确认 \\
\bottomrule
\end{tabularx}
\vfill
\callout{追问}{为什么“两个不相交的多数派”不可能同时存在？如果网络分区两边各有 2/5 个副本，能不能继续提交？}
\end{frame}

\begin{frame}{课堂推演：AI 模型别名切换}
\begin{itemize}
  \item v1 正在服务，v2 通过评测；
  \item 两个发布控制器同时尝试把 alias=prod 指向不同版本；
  \item 用共识日志、版本号和幂等 request ID 保证只有一个顺序；
  \item 推理服务收到事件后做滚动更新，失败则追加 rollback 命令，而不是覆盖历史。
\end{itemize}
\end{frame}

\sectionframe{七、Raft 完整时间线}{从选举、复制到提交，追踪一个命令如何成为所有副本的状态。}

\begin{frame}{Raft 例子：初始状态}
\begin{center}
\begin{tikzpicture}[>=Latex,node distance=12mm and 14mm]
\nodebox[minimum width=2.4cm]{ADSBlue}{a14}{A\\Follower t=4};
\nodebox[right=of a14,minimum width=2.4cm]{ADSOrange}{b14}{B\\Candidate t=5};
\nodebox[right=of b14,minimum width=2.4cm]{ADSGreen}{c14}{C\\Follower t=4};
\draw[\arrow,ADSOrange,thick] (b14)--node[above,font=\scriptsize]{RequestVote}(a14);
\draw[\arrow,ADSOrange,thick] (b14)--node[above,font=\scriptsize]{RequestVote}(c14);
\end{tikzpicture}
\end{center}
\begin{itemize}
  \item B 发现选举超时，递增 term=5 并投给自己；
  \item 收到多数票后成为 leader，先发送空 AppendEntries 建立权威；
  \item 旧 term=4 的 leader 如果恢复，必须接受 term=5 的更高任期。
\end{itemize}
\end{frame}

\begin{frame}{Raft 例子：复制命令}
\begin{enumerate}
  \item 客户端向 B 发送 SetAlias(v3)；
  \item B 将命令追加到本地日志 index=8、term=5；
  \item B 向 A、C 发送 AppendEntries(index=8,term=5)；
  \item A、C 校验 prevLogIndex/prevLogTerm 后追加；
  \item B 收到多数确认，但此时还要检查该 entry 是否属于当前 term。
\end{enumerate}
\end{frame}

\begin{frame}{Raft 例子：何时提交？}
\begin{center}
\begin{tikzpicture}[>=Latex]
\draw[ADSGray,thick] (0,0)--(11,0);
\foreach \x/\lab in {0/1,2/2,4/3,6/4,8/5,10/6}{\draw[ADSGray] (\x,.1)--(\x,-.1) node[below,font=\scriptsize]{\lab};}
\draw[ADSBlue,very thick] (6,.35)--(6,-.35) node[above,font=\scriptsize]{多数副本已有 index=4};
\node[font=\scriptsize,anchor=west] at (0,.65){term=5：index 4 = SetAlias(v3)};
\end{tikzpicture}
\end{center}
\begin{itemize}
  \item Leader 将 commitIndex 推进到 4；
  \item 每个副本按顺序 apply 到状态机，alias 才真正改变；
  \item 客户端收到成功响应应以 commit/应用确认作为语义边界，而非仅凭 leader 本地追加。
\end{itemize}
\end{frame}

\begin{frame}{旧 Leader 恢复：为什么不会双写？}
\begin{enumerate}
  \item A 恢复后仍带 term=4，发送旧日志；
  \item B/C 返回 term=5，A 更新 currentTerm 并退回 follower；
  \item 若 A 的日志与多数派不一致，AppendEntries 的 prevLog 检查失败；
  \item B 递减 nextIndex，找到最后一个匹配前缀；
  \item A 删除未提交冲突后缀，复制 term=5 的合法日志。
\end{enumerate}
\end{frame}

\begin{frame}{安全性证明的四个支点}
\begin{itemize}
  \item Election restriction：候选者必须拥有至少不落后的日志；
  \item Log matching：相同 index/term 的 entry 之前前缀完全相同；
  \item Leader completeness：已提交 entry 必在后续 leader 日志中；
  \item State machine safety：不同副本不会把不同命令应用到同一 index。
\end{itemize}
\vfill
\callout{证明方法}{先证明日志前缀不分叉，再证明 commitIndex 单调，最后推出状态机执行一致。}
\end{frame}

\begin{frame}{网络分区时两边能不能提交？}
\begin{tabularx}{\linewidth}{>{\bfseries}p{2.8cm}XX}
\toprule
3 副本 & 2/3 一侧可提交；1/3 一侧只能读旧状态或等待。 & 多数派相交。 \\
5 副本 & 3/5 一侧可提交；2/5 一侧不能提交新日志。 & 可容忍 2 个故障。 \\
双区域各半 & 两边都只有少数派。 & 安全但不可用，不能假装成功。 \\
\bottomrule
\end{tabularx}
\end{frame}

\begin{frame}{快照：日志太长之后怎么办？}
\begin{enumerate}
  \item 状态机在 index=10000 生成 snapshot；
  \item snapshot 包含应用状态、lastIncludedIndex 和 term；
  \item leader 可以向落后的副本发送 InstallSnapshot；
  \item 副本安装快照后，再接收 index>10000 的日志；
  \item 快照不是“随便保存一份状态”，必须与日志边界绑定。
\end{enumerate}
\end{frame}

\begin{frame}{成员变更：为什么不能直接替换节点？}
\begin{itemize}
  \item 从 3 节点直接切换到另一个 3 节点集合，可能出现两个不相交多数派；
  \item joint consensus 让旧配置和新配置同时参与提交；
  \item 只有 joint entry 提交后，才切换到新配置；
  \item 节点加入、移除和跨故障域迁移都需要同样的交集证明。
\end{itemize}
\end{frame}

\begin{frame}{模型注册中心的共识日志}
\begin{tabularx}{\linewidth}{>{\bfseries}p{2.8cm}XX}
\toprule
命令 & 日志语义 & 外部动作 \\
Publish(v3) & 记录模型 digest、指标、签名 & 写 registry 元数据 \\
Promote(v3) & 生产 alias 的唯一顺序 & 触发推理滚动发布 \\
Rollback(v2) & 追加补偿命令，不覆盖历史 & 逐步排空 v3 \\
Revoke(v3) & 撤销安全状态 & 拒绝新请求并清理缓存 \\
\bottomrule
\end{tabularx}
\end{frame}

\begin{frame}{共识日志与大对象的边界}
\begin{itemize}
  \item 共识日志只存模型 digest、URI、版本和策略，不存数百 GB 权重；
  \item 权重放在对象存储或内容寻址系统，发布前验证摘要与签名；
  \item 日志提交解决“哪个版本先成为生产版本”；
  \item 对象存储解决“权重在哪里、是否完整、如何复制”。
\end{itemize}
\end{frame}

\begin{frame}{课堂练习：提交还是拒绝？}
\begin{enumerate}
  \item 5 副本中 3 个收到 entry，但其中一个 entry term 比当前 leader 旧；
  \item leader 已知 commitIndex=7，客户端请求 index=8；
  \item 讨论：是否可以立即返回成功？还需要哪些确认？
  \item 延伸：如果外部模型 registry 写成功但共识提交失败，如何补偿？
\end{enumerate}
\vfill
\callout{答案关键词}{当前 term、日志匹配、commitIndex、幂等 request ID、outbox/fencing。}
\end{frame}

\begin{frame}{本讲总结}
\begin{enumerate}
  \item 共识要求 Termination、Agreement、Integrity、Validity；
  \item 同步模型允许用轮次解决 Crash，共识在纯异步模型中需要额外假设；
  \item 多数派交集、任期和日志匹配保证已提交历史不分叉；
  \item Raft 提供易理解的 leader/log/state machine 结构，Paxos 强调提案不变量；
  \item Kubernetes etcd、模型注册和发布系统都需要共识，但外部副作用仍需幂等与补偿。
\end{enumerate}
\vfill
\callout{下一讲衔接}{共识日志解决顺序，下一讲进入分布式存储、复制、分片和一致性读写。}
\end{frame}

\begin{frame}{课后题}
\begin{enumerate}
  \item 说明 FLP 直觉中“慢”和“崩溃”为什么不可区分；
  \item 画出 Raft 一次 leader 选举和日志提交的时序；
  \item 证明两个多数派一定相交；
  \item 设计模型注册中心的共识日志命令和回滚命令。
\end{enumerate}
\end{frame}

\begin{frame}{Raft/Paxos 结论的边界}
\small
\begin{itemize}
  \item Raft 中，leader 不能仅因为某条日志被多数副本保存就无条件提交；提交规则还要结合当前任期，旧任期日志通常要通过提交一条当前任期日志间接确认。
  \item Paxos 中 accepted、chosen、learned 是三个不同状态：被某个 acceptor 接受不等于已经被多数派选定，也不等于所有学习者都已应用。
  \item 网络分区时，少数派可以继续服务部分只读或缓存请求，但不能安全地产生新的多数派提交；“可访问”不等于“可线性化写入”。
  \item 共识日志只解决复制状态机的顺序和决定；模型仓、对象存储、GPU 调度等外部副作用仍需幂等、补偿和 fence。
\end{itemize}
\vfill
\callout{证明习惯}{看到“已提交”“已决定”“已应用”时，分别写出证据：多数派、日志匹配、学习/应用位置，而不要只看网络上已有几份副本。}
\end{frame}

\end{document}
